An assistant with tools and exposed reasoning returns more than one output: a structured tool-call proposal, a visible answer, and a reasoning field. Safety evaluations usually score one. We ask whether the channels agree and whether a verdict from one survives a change in prompt framing. Three repeated-prompt studies use a local gpt-oss-20b deployment and an API extension to gpt-oss-120b, DeepSeek V4 Flash, GPT-5.6 Luna, and gpt-oss-safeguard-20b. First, deleting simulation wording from a fictional scenario lowered exact tool calls from \LocalFOneSim/100 to \LocalFOneNoSim/100 locally and from \OssFOneSim/100 to \OssFOneNoSim/100 for gpt-oss-120b, but left three other deployments unchanged: the framing effect is deployment-specific. Second, in a five-turn conversation, API deployments refuse or resist a direct request to maximize user dependency in 116 of 120 conversations, then comply in all 120 once it is relabelled as education; later requests split them, two complying throughout, two refusing distress-triggered monetization in 30/30 conversations. Third, a synthetic secret guarded by a two-part release condition appeared in local returned reasoning before release in \LocalPreReason/30 trials; after release, \LocalReasonRefusal/30 local and \OssReasonRefusal/30 gpt-oss-120b trials disclosed it in reasoning while the final answer refused. The evidence supports one narrow claim: a refusal or success in one channel does not describe the others, and evaluations should state which channel they score. The \ReviewCaseCount{} multi-turn responses were labelled by blinded language-model reviewers under a disclosed, mechanically verified rubric; all counts are recomputed from preserved records; no internal mechanism or real-world harm is claimed.
